\documentclass{article}

% Language setting
% Replace `english' with e.g. `spanish' to change the document language
\usepackage[czech]{babel}
\usepackage{amsthm,thmtools,xcolor,amsmath,amssymb, mathtools}

% Set page size and margins
% Replace `letterpaper' with `a4paper' for UK/EU standard size
\usepackage[a4paper,top=2cm,bottom=2cm,left=3cm,right=3cm,marginparwidth=1.75cm]{geometry}

% Useful packages
\usepackage{amsmath}
\usepackage{enumerate}
\usepackage{graphicx}
\usepackage{tabularx}
\usepackage[colorlinks=true, allcolors=blue]{hyperref}
\graphicspath{images/}

%\usepackage{tikzit}
%\input{default.tikzstyles}

\title{DÚ Lineární algebra -- Sada 6}
\author{Jan Romanovský}

\begin{document}
\maketitle

\textbf{(6.2)}
    Vektory $(2, 1, 0, 1)^T$ a $(1,0,1,2)$ jsou přímo báze (ozn. $B$) prostoru $\textbf{V}$, vyjádřeme si obrazy vektorů v této bázi:
    $$\left(\begin{tabular}{ c c | c c }
        2 & 1 & 4 & 2\\
        1 & 0 & 1 & 2\\
        0 & 1 & 2 & -2\\
        1 & 2 & 5 & -2
    \end{tabular}\right)\sim
    \left(\begin{tabular}{ c c | c c }
        1 & 0 & 1 & 2\\
        0 & 1 & 2 & -2\\
        0 & 0 & 0 & 0 \\
        0 & 0 & 0 & 0
    \end{tabular}\right)$$
    Matice operátoru v bázi $B$ je tedy:
    $$\begin{pmatrix}
        1 & 2 \\
        2 & -2
    \end{pmatrix}$$
    Když spočítame vlastní čísla a vlastní vektory:
    $$p_b(\lambda) = (1-\lambda)(-2-\lambda)-4, \lambda_1 = -3, \lambda_2 = 2, \textbf{v}_{\lambda_1} \in \text{span}\left\{\begin{tabular}{c}1 \\ -2  \end{tabular}\right\}, \textbf{v}_{\lambda_2} \in \text{span}\left\{\begin{tabular}{c} 2 \\ 1 \end{tabular}\right\}$$
        To nám dává novou bázi (ozn. $C$) z vlastních vektorů v níž je matice operátoru diagonální, vyberme $(1, -2)^T, (2, 1)^T$, jejich matici přechodu k bázi $B$ vidíme zřejmě; když tedy dáme vše dohromady (ozn. matici operátoru $f$ v kanonické bázi jako $A$)(první a poslední matice nejsou přesně matice přechodu, protože $B$ a $K$ mají jinou dimenzi, ale je to analogické, poslední matice je lib. levý inverz první matice):
       \begin{align*}
           A &= [\text{id}]_K^B[\text{id}]_B^C\begin{pmatrix}
   -3 & 0 \\ 0 & 2
        \end{pmatrix}[\text{id}]_C^B[\text{id}]_B^K \\&= [\text{id}]_K^B[\text{id}]_B^C\begin{pmatrix}
   -3 & 0 \\ 0 & 2
        \end{pmatrix}\left\{[\text{id}]_B^C\right\}^{-1}\left\{[\text{id}]_K^B\right\}^{-1} \\&= \begin{pmatrix}
    2 & 1 \\ 1 & 0 \\ 0 & 1 \\ 1 & 2
        \end{pmatrix}\begin{pmatrix}
    1 & 2 \\ -2 & 1
        \end{pmatrix}\begin{pmatrix}
   -3 & 0 \\ 0 & 2
 \end{pmatrix}\left(\begin{array}{c c}
   \frac{1}{5} & -\frac{2}{5} \vspace{0.25em}\\ \frac{2}{5} & \frac{1}{5}
 \end{array}\right)\begin{pmatrix}
            0 & 1 & 0 & 0 \\ 1 & -2 & 0 & 0
        \end{pmatrix}
       \end{align*}
        Mocnění operátoru už je nyní jednoduché:
        \begin{align*}
            f^n\left((5, 1, 3, 7)^T\right) &= \begin{pmatrix}
    2 & 1 \\ 1 & 0 \\ 0 & 1 \\ 1 & 2
        \end{pmatrix}\begin{pmatrix}
    1 & 2 \\ -2 & 1
        \end{pmatrix}\begin{pmatrix}
   -3 & 0 \\ 0 & 2
 \end{pmatrix}^n\left(\begin{array}{c c}
   \frac{1}{5} & -\frac{2}{5} \vspace{0.25em} \\ \frac{2}{5} & \frac{1}{5}
        \end{array}\right)\begin{pmatrix}
            0 & 1 & 0 & 0 \\ 1 & -2 & 0 & 0
        \end{pmatrix}\begin{pmatrix}
            5 \\ 1 \\ 3 \\ 7
        \end{pmatrix}
            \\ &= \begin{pmatrix}
    2 & 1 \\ 1 & 0 \\ 0 & 1 \\ 1 & 2
        \end{pmatrix}\begin{pmatrix}
    1 & 2 \\ -2 & 1
        \end{pmatrix}\begin{pmatrix}
            -3^n & 0 \\ 0 & 2^n
        \end{pmatrix}\left(\begin{array}{c c}
   \frac{1}{5} & -\frac{2}{5} \vspace{0.25em}\\ \frac{2}{5} & \frac{1}{5}
        \end{array}\right)\begin{pmatrix}
            0 & 1 & 0 & 0 \\ 1 & -2 & 0 & 0
        \end{pmatrix}
            \\ &= \begin{pmatrix}
2^n & 0 & 0 & 0 \\
\frac{1}{5} (-2) (-3)^n+\frac{2^{n+1}}{5} & -2 \left(\frac{1}{5} (-2) (-3)^n+\frac{2^{n+1}}{5}\right)+\frac{(-3)^n}{5}+\frac{2^{n+2}}{5} & 0 & 0 \\
\frac{4 (-3)^n}{5}+\frac{2^n}{5} & -2 \left(\frac{4 (-3)^n}{5}+\frac{2^n}{5}\right)+\frac{1}{5} (-2) (-3)^n+\frac{2^{n+1}}{5} & 0 & 0 \\
\frac{1}{5} (-2) (-3)^{n+1}+\frac{2^{n+2}}{5} & -2 \left(\frac{1}{5} (-2) (-3)^{n+1}+\frac{2^{n+2}}{5}\right)+\frac{1}{5} (-3)^{n+1}+\frac{2^{n+3}}{5} & 0 & 0
          \end{pmatrix}\begin{pmatrix}
            5 \\ 1 \\ 3 \\ 7
        \end{pmatrix}
            \\ &= \begin{pmatrix}
 5\cdot 2^n \\
 3 \left(\frac{1}{5} (-2) (-3)^n+\frac{2^{n+1}}{5}\right)+\frac{(-3)^n}{5}+\frac{2^{n+2}}{5} \\
 3 \left(\frac{4 (-3)^n}{5}+\frac{2^n}{5}\right)+\frac{1}{5} (-2) (-3)^n+\frac{2^{n+1}}{5} \\
 3 \left(\frac{1}{5} (-2) (-3)^{n+1}+\frac{2^{n+2}}{5}\right)+\frac{1}{5} (-3)^{n+1}+\frac{2^{n+3}}{5}
        \end{pmatrix}
        \end{align*}
\end{document}
